%\documentclass[preprint,tighten,prl]{problemDB}
\documentclass[]{revtex4}
%\documentclass[]{book}
%\newcommand{\sm}[1]{{\[\]}}
\newcommand{\sm}[1]{{\[#1\]}}
%\newcommand{\st}[1]{#1}
\newcommand{\st}[1]{}

\usepackage{amsthm}
\usepackage{amsmath}
\usepackage{hyperref}
\usepackage{graphicx} % Include figure files
\usepackage{bm}
\usepackage{float}
\usepackage[usenames,dvipsnames,svgnames,table]{xcolor}
\usepackage{siunitx}
\usepackage{ulem}
\usepackage{url}
\usepackage{enumerate}

\listfiles
\renewcommand{\vec}{\mathbf}
\renewcommand{\[}{\begin{equation}}
\renewcommand{\]}{\end{equation}}
%\newcommand{\pzi}{\pause\item }
\newcommand{\bfr}{\begin{frame}}
\newcommand{\efr}{\end{frame}}
\newcommand{\be}{\begin{enumerate}}
\newcommand{\ee}{\end{enumerate}}
\newcommand{\bi}{\begin{itemize}}
\newcommand{\ei}{\end{itemize}}
\newcommand{\bsi}{\begin{itemize}}
\newcommand{\esi}{\end{itemize}}
\newcommand{\cork}[1]{}    %ideas for corks will pop up when making the worksheets.  Write them down embedded in the text and we can parse them out later

%['0B3C5D', '062F4F', '328CC1', 'D9B310', '984B43', 'B82601', '57652A', '76323F', '626E60', 'AB987A', 'C09F80', 'b0b0b0']
\definecolor{bg0}{HTML}{0B3C5D}
\definecolor{bg5}{HTML}{B82601}
\newcommand{\sjb}[1]{{\bf [sjb: #1]}}
\newcommand{\zz}[1]{{\color{blue} [zz: #1]}}
\newcommand{\qz}[1]{{\color{Purple} [qz: #1]}}
%\textcolor{bg0}{#1}}


\begin{document}

\title{Physics of Crystalline Materials\\
MSAE 4200\\
Worksheet: Group Theory}

\maketitle

\vskip 24pt
{Names and Roles of team members:}

\vskip 24 pt
Goal: The goals of this spreadsheet are:
\begin{itemize}
\item To develop the basic tools for understanding the important fluctuation-dissipation theorem
\item To understand the fluctuation-dissipation theorem
\end{itemize}

\cork{}
\vskip 12pt



%\vskip 12pt
Start this worksheet in class in your groups and discuss it.  There will also be some class discussion to get things off to a good start.
This worksheet will be finished in groups for  homework.

%\section*{What is this worksheet about?}
\section*{OK, here we go:}
Usual preamble:
\begin{enumerate}
  \item Figure out your group
  \item Figure out who is team leader
  \item Introduce yourselves to each other
  \item Figure out who is scribe for this worksheet
  \item Discuss and write down the team goal
  \item We will work on this worksheet in groups in class but you will finish up by yourself for homework.
\end{enumerate}

\section*{Review of where we are}
\bi
\item We learned that if we have linear problems there are some rather general results that are true, in particular, linear problem + causality
mean that we can derive a definite relationship between the real and imaginary parts of the susceptibility.
\item This means that if we can measure either the real-part, or the imaginary part, over a sufficiently wide range of frequencies,
we can determine the one we don't know frm the other using the Kramers-Kronig relations.
\item The imaginary part is gives the dissipation in the system
\ei

\section*{}
\begin{enumerate}[(a)]
  \item 
\end{enumerate}
\end{document}
